\documentclass[a4paper,12pt]{article}
\usepackage[utf8]{inputenc}
\usepackage[T1]{fontenc}
\usepackage[brazil]{babel}
\usepackage{geometry}
\usepackage{listings}
\usepackage{xcolor}
\usepackage{longtable}
\usepackage{changepage}
\geometry{margin=2.5cm}

\title{DPP Compiler \\ GLC da Linguagem ``Dhara''}
\author{Pedro Schneider e Gabriel Santana Dias}
\date{Abril 2026}
\lstset{
    basicstyle=\ttfamily\small,
    keywordstyle=\color{blue},
    commentstyle=\color{gray},
    stringstyle=\color{orange},
    breaklines=true,
    frame=single,
    showspaces=false,
    showstringspaces=false
}

\begin{document}
\maketitle
\section{Introdução}
Este projeto implementa uma linguagem fictícia chamada \textbf{Dhara} (DPP), cujas palavras reservadas são baseadas em títulos de músicas.
\section{Palavras-chave da Linguagem}
\begin{longtable}{|l|c|l|}
\hline
\textbf{Função} & \textbf{DPP} & \textbf{Tradução} \\
\hline
Início do programa & style & --- \\
Fim do programa & borderline & --- \\
Comentário de linha & $\sim\sim$ & // (em C) ou \# (em Python) \\
Tipo inteiro & space & inteiro \\
Tipo decimal & lithium & decimal \\
Tipo string & judas & texto \\
Função & PREY & function \\
Retorno & HOMETOWN & return \\
Saída & catapult & output \\
Entrada & pleaser & input \\
If & houdini & se \\
Else & more & senao \\
While & problems & enquanto \\
For & bloomfield & para \\
Do...While & not...ok & faça enquanto \\
Pare & loser & pare \\
\hline
\end{longtable}

\newpage
\section{Gramática Livre de Contexto (GLC)}
\begin{adjustwidth}{-2cm}{-2cm}
\begin{lstlisting}[basicstyle=\ttfamily\footnotesize]
palavra    -> [a-zA-Z]+[a-zA-Z]*
num        -> [0-9]+

id         -> palavra+num*
int        -> num+
float      -> num+ '.' num+
string     -> " palavra* "

op_logic   -> && | ||
op_comp    -> == | != | >= | <= | > | <
op_arit    ->  + |  - |  * |  / | %

tipo       -> `space` | `lithium` | `judas`
atribuicao -> tipo id = expressao ; | id = expressao ;
declaracao -> tipo id ;

comentario -> `~~` palavra* `\n`

expressao  -> termo expressao'
expressao' -> op_arit termo expressao' | e
termo      -> ( expressao ) | int | float | string | id | chamada_funcao

condicao   -> termo condicao'
condicao'  -> (op_logic | op_comp) termo condicao' | e

saida             -> `catapult` [ expressao ] ;

entrada           -> `pleaser` [ entrada' ] ;
entrada'          -> " ponteiros " , identificadores
ponteiros         -> formatadores ponteiros'
ponteiros'        -> , ponteiros | e
identificadores   -> id identificadores'
identificadores'  -> , identificadores | e
formatadores      -> %d | %f | %s

declaracao_funcao       -> tipo `PREY` id [ parametros_declaracao ] { comando* `HOMETOWN` expressao? }
parametros_declaracao   -> tipo id parametros_declaracao'
parametros_declaracao'  -> , tipo id parametros_declaracao' | e

chamada_funcao          -> id [ parametros_chamada ] ;
parametros_chamada      -> expressao parametros_chamada'
parametros_chamada'     -> , expressao parametros_chamada' | e

if       -> `houdini` ( condicao ) { comando* } else
else     -> `more` { comando* } | e
while    -> `problems` ( condicao ) { comando* }
do_while -> `not...ok` { comando* } while ( condicao ) ;
for      -> `bloomfield` ( atribuicao condicao atribuicao ) { comando* }

comando  -> comentario | atribuicao | entrada | saida | if | while | do_while | for | chamada_funcao
main     -> declaracao_funcao* `style` codigo `borderline`
codigo   -> comando* codigo* | e
\end{lstlisting}
\end{adjustwidth}

\lstdefinelanguage{Dhara}{
    % keywords by group — each group can have its own color
    morekeywords=[1]{style, borderline, PREY, HOMETOWN},
    morekeywords=[2]{space, lithium, judas},
    morekeywords=[3]{houdini, more, problems, bloomfield, loser},
    morekeywords=[4]{catapult, pleaser},
    sensitive=true,        % case-sensitive keywords
    morestring=[b]",       % strings delimited by "
    morecomment=[l]{~~},   % line comments starting with ~~
}

\lstdefinestyle{DharaStyle}{
    language=Dhara,
    basicstyle=\ttfamily\small,
    keywordstyle=[1]\color{violet}\bfseries,   % style, borderline, PREY, HOMETOWN
    keywordstyle=[2]\color{blue}\bfseries,     % space, lithium, judas
    keywordstyle=[3]\color{orange}\bfseries,   % houdini, more, problems, bloomfield
    keywordstyle=[4]\color{teal}\bfseries,     % catapult, pleaser
    stringstyle=\color{red},
    commentstyle=\color{gray}\itshape,
    showstringspaces=false,
    frame=single,
    breaklines=true,
}

\section{Exemplos de Código}
\subsection{Declaração de Variável}
\begin{lstlisting}[style=DharaStyle]
space x;        ~~ int x;
lithium y;      ~~ float y;
judas texto;    ~~ string texto;
\end{lstlisting}

\subsection{Estrutura If-Else}
\begin{lstlisting}[style=DharaStyle]
houdini(x > 10) {
    catapult["Maior que 10"];
} more {
    catapult["Menor ou igual a 10"];
}
\end{lstlisting}

\subsection{Laço de Repetição}
\begin{lstlisting}[style=DharaStyle]
bloomfield(x=0; x<10; x++) {
    catapult[x];
}
\end{lstlisting}

\subsection{Declarar Função}
\begin{lstlisting}[style=DharaStyle]
space PREY funcao_que_retorna_inteiro[space x, space y] {
    houdini(y != 0){
        space resultado = x/y;
    } more {
        HOMETOWN;
    }
    HOMETOWN resultado;
}
\end{lstlisting}

\newpage
\section{Exemplo Completo}
\subsection{Código em Dhara}
\begin{lstlisting}[style=DharaStyle]
lithium PREY multiplicar[lithium a, lithium b] {
    ~~ multiplica dois numeros
    HOMETOWN a*b;
}

style
space x;
space q;
lithium y;

bloomfield(x=0; x<=10; x++) {
    catapult[x];
}

houdini((x == 10) && (x != 0)) {
    catapult["Digite um numero inteiro e um decimal: "];
    pleaser["%d,%f", q, y];
    catapult["Voce digitou: " + q + " e " + y];
}

houdini(y < 10.5) {
    problems(y < 11) {
        y++;
    }
} more {
    lithium res = multiplicar[q, y];
    catapult["Q*Y e igual a: " + res];
}
borderline
\end{lstlisting}

\newpage
\subsection{Código Equivalente em C}
\begin{lstlisting}[language=C]
float multiplicar(float a, float b) {
    // multiplica dois numeros
    return a*b;
}

int main(){
  int x;
  int q;
  float y;

  for(x=0; x<=10; x++){
    printf("%d", x);
  }

  if((x==10) && (x!=0)){
    printf("Digite um numero inteiro e um decimal: ");
    scanf("%d %f", &q, &y);
    printf("Voce digitou: %d e %f", q, y);
  }

  if(y < 10.5){
    while(y < 11){
        y++;
    }
  } else{
    float res = multiplicar(q, y);
    printf("Q*Y e igual a: %f", res);
  }
}
\end{lstlisting}
\end{document}